% !Mode:: "TeX:UTF-8"
%%% Local Variables:
%%% mode: latex
%%% TeX-master: t
%%% End:
\renewcommand\theequation{A.\arabic{equation}}
\renewcommand\thefigure{A.\arabic{figure}}
\renewcommand\thetable{A.\arabic{table}}
\setcounter{equation}{0} 
\setcounter{figure}{0}
\setcounter{table}{0}
%%%============================================

\chapter*{附\ \ 录}\label{Appendix}
  \addtocontents{toc}{\vspace{-9pt}}% 在自动生成的目录中减少 9pt的垂直间距
\addcontentsline{toc}{chapter}{附录{\cdotfill}} 
\vspace{0.68cm}

%%%%%   附录一
%\section*{附录一 \quad 逻辑推理素养水平划分}\label{fulu2}
%\addcontentsline{toc}{section}{附录一 \quad 逻辑推理素养水平划分}
%\renewcommand\theequation{1.\arabic{equation}}
%\renewcommand\thefigure{1.-\arabic{figure}}
%\renewcommand\thetable{1.\arabic{table}}
%\setcounter{equation}{0}
%\setcounter{figure}{0}
%\setcounter{table}{0}

小二号黑体居中，附录二字与题名间空一个汉字符位，然后隔行编排正文主体内容。 

\section*{附录一 \quad 高中生逻辑推理素养前测试卷}\label{fulu2}
\addcontentsline{toc}{section}{附录一\quad 高中生逻辑推理素养前测试卷}
\renewcommand\theequation{1.\arabic{equation}}
\renewcommand\thefigure{1.-\arabic{figure}}
\renewcommand\thetable{1.\arabic{table}}
\setcounter{equation}{0}
\setcounter{figure}{0}
\setcounter{table}{0}

%\subsection*{A.1\quad 高中生逻辑推理素养初测试卷} % 附录的二级标题不是必须的
%\addcontentsline{toc}{subsection}{A.1\quad 高中生逻辑推理素养初测试卷}
亲爱的各位同学, 你好！非常感谢您参加本次调查研究, 你将有 50 分钟来完成本次测
试卷. 本次试卷仅作为相关研究资料, 作答记录不会记录到你的档案中, 请根据实际情况作
答. 你的作答对我们的研究十分重要, 再次感谢你的支持与合作.\\
\\
第一部分: 以下为基础信息调查, 请如实填写\\
请问你的性别是(\quad )\\
A. 男\qquad  \qquad \qquad B. 女\\
请问你上次月考的数学成绩是(\quad  )\\
A. 120-150 \qquad  \qquad \qquad B. 100-119\qquad \qquad \qquad C. 80-99\qquad \qquad \qquad D. 60-79  \\
E.其它\\
\\
第二部分:测试题\\
1. 天文学家经研究认为“地球和火星在太阳系中有某些相似的性质, 而地球上有生命, 因
此推测火星上也有生命存在”, 这里运用到了什么推理方法(\quad )\\
A. 归纳推理\qquad  \qquad \quad B.类比推理 \; \quad  \quad \qquad C.演绎推理 \ \quad  \quad \qquad D.反证法\\
\\
\\
2. 三角函数都是周期函数(大前提), $tan\alpha$是三角函数(小前提), 因此$tan\alpha$ 是周期函数(结论). 这里用到的推理方法是(\quad  ), 并写出一个和之相似的命题\\
A. 归纳推理\qquad  \qquad \quad B. 类比推理 \; \quad  \quad \qquad C. 演绎推理 \ \quad  \quad \qquad D. 合情推理\\
\\
\\
 
 
 
 \newpage
\section*{附录二 \quad 高中生逻辑推理素养后测试卷}\label{fulu2}
\addcontentsline{toc}{section}{附录二 \quad 高中生逻辑推理素养后测试卷}
\renewcommand\theequation{1.\arabic{equation}}
\renewcommand\thefigure{1.-\arabic{figure}}
\renewcommand\thetable{1.\arabic{table}}
\setcounter{equation}{0}
\setcounter{figure}{0}
\setcounter{table}{0}
%\subsection*{A.2\quad 高中生逻辑推理素养后测试卷}  % 附录的二级标题不是必须的
%\addcontentsline{toc}{subsection}{A.2 \quad  高中生逻辑推理素养后测试卷}
亲爱的各位同学, 你好！ 非常感谢您参加本次调查研究, 你将有 50 分钟来完成本次测
试卷. 本次试卷仅作为相关研究资料, 作答记录不会记录到你的档案中, 请根据实际情况作
答. 你的作答对我们的研究十分重要, 再次感谢你的支持与合作.\\
\\
班级:    \underline{ \qquad  \qquad \qquad }        \qquad  \qquad \qquad             性别:  \underline{ \qquad  \qquad \qquad }    \\
\\  
1. 观察下列图片,并阅读下面的文字表述, 则当6条直线相交时, 最多有(\quad)个交点.\\
 
2. 从一个点引出三条不在同一平面内的射线, 用一个平面截这三条射线得到的图形是三棱锥. 从一个点引出四条射线, 其中任意三条射线不在同一平面, 用一个平面截这四条射线, 所得到的图形就是四棱锥. 请同学们仿照三棱锥和四棱锥的定义, 尝试给出五棱锥的定义.\\
 

%%%%%   附录 B
%\chapter*{}

\section*{附录三 \quad 教师访谈提纲}\label{fulu2}
\addcontentsline{toc}{section}{附录三 \quad 教师访谈提纲}
\renewcommand\theequation{1.\arabic{equation}}
\renewcommand\thefigure{1.-\arabic{figure}}
\renewcommand\thetable{1.\arabic{table}}
\setcounter{equation}{0}
\setcounter{figure}{0}
\setcounter{table}{0}
 
 
% \subsection*{B.1\quad 教师访谈提纲}  % 附录的二级标题不是必须的
%\addcontentsline{toc}{subsection}{B.1\quad 教师访谈提纲}
尊敬的老师:

您好！

为更好地了解一线高中教师对逻辑推理素养的理解情况, 特进行此次访谈, 感谢您百忙之中抽出时间参加. 访谈结果仅供本人论文研究使用, 您可以根据您的实际情况进行作答, 再次感谢您的参与.\\
\\
1.	谈谈您对逻辑推理素养的认识.\\
2.	您觉得目前高中生的逻辑推理素养水平如何?\\
3.	您认为影响学生逻辑推理素养的因素有哪些?\\
4.	您认为在日常教学中该如何培养学生的逻辑推理素养? 在``圆锥曲线的方程"这一章节中有何针对性培养措施?\\
5.	谈谈您对《高考评价体系》的认识.\\
6.	您认为《高考评价体系》中提出的``一核"、``四层"与``四翼"对学生的逻辑推理素养培养有指导意义吗？



 
